\begin{figure}[tb]
\centering
\begin{tikzpicture}[>=Stealth,font=\small]
 \begin{scope}[xshift=-3.5cm,yshift=0.4cm]
  \coordinate (u) at (0,0);
  \coordinate (v) at (1.45,0);
  \coordinate (w) at (2.9,0);
  \draw[thick,->] (u)--(v);
  \draw[thick] (v)--(w);
  \draw[thick] (u) .. controls (-0.85,1.05) and (0.85,1.05) .. (u);
  \draw[thick] (w) .. controls (2.05,1.05) and (3.75,1.05) .. (w);
  \foreach \x in {u,v,w}{\fill (\x) circle (1.8pt);}
  \node[below=4pt] at (u) {$u$};
  \node[below=4pt] at (v) {$v$};
  \node[below=4pt] at (w) {$w$};
  \node at (1.45,-0.95) {$(\deg(u),\deg(v),\deg(w))=(3,2,3)$};
  \node at (1.45,1.45) {map: four edges};
 \end{scope}
 \begin{scope}[xshift=2.8cm,yshift=1.55cm,
  block/.style={circle,draw,fill=black!10,inner sep=2pt,minimum size=15pt},
  leaf/.style={circle,draw,inner sep=1.5pt,minimum size=13pt}]
  \node[block] (r) at (0,0) {$u$};
  \node[block] (a) at (-1.15,-0.85) {$u$};
  \node[block] (b) at (1.15,-0.85) {$v$};
  \node[leaf] (a1) at (-1.55,-1.65) {$u$};
  \node[leaf] (a2) at (-0.75,-1.65) {$u$};
  \node[leaf] (b1) at (0.7,-1.65) {$v$};
  \node[block] (b2) at (1.6,-1.65) {$w$};
  \node[leaf] (c1) at (1.2,-2.45) {$w$};
  \node[leaf] (c2) at (2.0,-2.45) {$w$};
  \foreach \x/\y in {r/a,r/b,a/a1,a/a2,b/b1,b/b2,b2/c1,b2/c2}
   {\draw (\x)--(\y);}
  \node[above=2pt] at (r.north) {link};
  \node[left=2pt] at (a.west) {loop};
  \node[right=2pt] at (b.east) {link};
  \node[right=2pt] at (b2.east) {loop};
  \node at (0.1,-3.02) {block tree: nine nodes};
 \end{scope}
\end{tikzpicture}
\caption{Corner substitution and pole multiplicities. Shaded nodes carry
one-edge blocks; white nodes are trivial inserts. Node labels give
their poles in the map. The eight nonroot nodes have pole counts
$3,2,3$, exactly the vertex degrees. The arrow marks the root dart.}
\label{fig:corners}
\end{figure}
